\chapter{Provability, the Second Incompleteness Theorem, and Truth}
\label{ch:second-incompleteness}

\section{Derivability conditions}
\label{sec:second-incompleteness-01}

Fix a computably enumerable theory \(T\supseteq\mathsf Q\) and a
specified arithmetical provability formula, denoted here by
\(\mathsf{Prov}_T(x)\). Its proof-existence meaning is
\(T\derives\sigma\) if and only if
\(\mathcal N\entails\mathsf{Prov}_T
(\overline{\ulcorner\sigma\urcorner})\), for arithmetic sentences
\(\sigma\). The certified construction of
\cref{def:fixed-proof-predicate} supplies one such presentation.
The theorem statements below require the stated internal derivability
conditions on the particular formula being used.

For the standard \(\PA\) applications in this chapter, make the explicit
local selection
\[
 \mathsf{Prov}_{\PA}(x)\defeq\mathsf{Prv}^{19}_{\PA}(x)
     =\exists p\,\Gamma_{d_{19}}(p,x,\bar1)
\]
from \cref{def:p05-selected-direct-pa-graph}.
Its proof-existence meaning is given by
\cref{thm:p05-selected-direct-pa-graph}. This is the direct tag-19
presentation; the earlier certified \(\mathsf{Proof}_{\PA}\) codes
and an arbitrarily selected representative of them remain distinct.
The parameter \(\mathsf{Prov}_T\) in each use of the following theorems
includes this formula choice, not just the theory name.
To reduce clutter, write
\(\mathsf{Prov}_T(\ulcorner\varphi\urcorner)\) for
\(\mathsf{Prov}_T(\overline{\ulcorner\varphi\urcorner})\).

\begin{definition}[Hilbert--Bernays--L\"ob derivability conditions]
The provability predicate satisfies the \emph{derivability conditions} when,
for all sentences \(\varphi\) and \(\psi\):
\begin{enumerate}
  \item if \(T\derives\varphi\), then
        \(T\derives\mathsf{Prov}_T(\ulcorner\varphi\urcorner)\);
  \item
  \[
    T\derives
    \mathsf{Prov}_T(\ulcorner\varphi\to\psi\urcorner)
    \to
    \bigl(\mathsf{Prov}_T(\ulcorner\varphi\urcorner)
    \to\mathsf{Prov}_T(\ulcorner\psi\urcorner)\bigr);
  \]
  \item
  \[
    T\derives
    \mathsf{Prov}_T(\ulcorner\varphi\urcorner)
    \to
    \mathsf{Prov}_T
    (\ulcorner\mathsf{Prov}_T(\ulcorner\varphi\urcorner)\urcorner).
  \]
\end{enumerate}
The first clause is a metatheoretic implication; the latter two are theorem
schemes internal to \(T\).  The three clauses are cited below as D1, D2,
and D3.
\end{definition}

For the selected \(\mathsf{Prv}^{19}_{\PA}\),
\cref{thm:p05-direct-pa-internal-mp} proves D1 and the full internal D2
scheme through its specified proof-joining description, and
\cref{thm:p05-direct-pa-d3} proves D3 through the original-sentence
evidence compiler. All three conditions therefore hold for the same
literal \(\exists p\,\Gamma_{d_{19}}(p,x,\bar1)\).
For an arbitrary computably enumerable \(T\supseteq\mathsf Q\), the
results below retain the three conditions as hypotheses on its chosen
formula \cite{HilbertBernays1939,Lob1955}.

\begin{definition}[Formal consistency sentence]
\label{def:formal-consistency-sentence}
Fix \(\bottom_A\) as the false arithmetic sentence \(0=S0\).  Define
\[
  \mathsf{Con}(T)\defeq
  \neg\mathsf{Prov}_T(\overline{\ulcorner\bottom_A\urcorner}).
\]
This formula expresses consistency relative to the fixed numbering and proof
predicate.
With the selected standard \(\PA\) convention in this chapter it is
specifically
\(\neg\mathsf{Prv}^{19}_{\PA}
(\overline{\ulcorner0=S0\urcorner})\).
\end{definition}

\begin{warning}[External and internal consistency]
The metalinguistic claim that \(T\) is consistent is not literally the formula
\(\mathsf{Con}(T)\).  Correctness of the chosen proof predicate connects them
in the standard structure \(\mathcal N\); the second incompleteness
theorem concerns whether
\(T\) proves this particular internal sentence.
\end{warning}
